% !TeX root = ../../../proposal.tex
% Front matter: title page, abstracts, TOC, list of vocab

% \pdfstringdefDisableCommands{
%   \let\MakeUppercase\relax
% }

% \begin{center}
%   \includegraphics[width=2.8cm]{figures/common/logo02.jpg}
% \end{center}
% \vspace*{-1cm}

% \maketitlepage
% \makesignaturepage

% %%%%%%%%%%%%%%%%%%%%%% English abstract %%%%%%%%%%%%%%%%%%%%%%%%
% \abstract
% [Write your English abstract here.]

% \justifying
% \begin{tabularx}{\textwidth}{@{}l>{\RaggedRight\arraybackslash}X@{}}
%   \textbf{Keywords}: & Keyword1 / Keyword2 / Keyword3
% \end{tabularx}
% \vspace{1mm}
% \endabstract

% %%%%%%%%%%%%%%%%%%%%%%%% Thai abstract %%%%%%%%%%%%%%%%%%%%%%%%%

% \thaiabstract
% \justifying
% [เขียนบทคัดย่อภาษาไทยของคุณที่นี่]

% \justifying
% \begin{tabularx}{\textwidth}{@{}l>{\RaggedRight\arraybackslash}X@{}}
%   & \\
%   \textbf{คำสำคัญ}: & คำสำคัญ1 / คำสำคัญ2 / คำสำคัญ3
% \end{tabularx}
% \vspace{1mm}
% \endabstract

% %%%%%%%%%%%%%%%%%%%%%%%%% Acknowledgments %%%%%%%%%%%%%%%%%%%%%%
% \preface
% [เขียนกิตติกรรมประกาศ/คำขอบคุณของคุณที่นี่]

% %%%%%%%%% Table of Contents, List of Figures/Tables %%%%%%%%%%%%
% %: The three commands below automatically generate the table
% %: of content, list of tables, list of figures, and list of listings

% %: Appendices or Appendix
% %: If you have more than one appendix, use `Appendices`
% \def\appendixnames{Appendices}

% \tableofcontents
% \listoftables
% \listoffigures

% %%%%%%%%%%%%%%%%%%%%%%%% List of codes %%%%%%%%%%%%%%%%%%%%%%%%%
% %: Command below automatically will generate the table
% %: for listing 'lst' package listings

% \listofprograms

%%%%%%%%%%%%%%%%%%%%%%%% List of symbols %%%%%%%%%%%%%%%%%%%%%%%
%: You have to add this manually...

% \listofsymbols
% \newline\justifying{
% \begin{tabular}{@{}p{0.07\textwidth}p{0.7\textwidth}p{0.1\textwidth}}
% \textbf{สัญลักษณ์} & & \textbf{หน่วย} \\[0.2cm]
%     $\alpha$ & Test variable\hfill & m$^2$ \\
%     $\lambda$ & Interarival rate\hfill &  jobs/second\\
%     $\mu$ & Service rate\hfill & jobs/second\\
% \end{tabular}
% }

%\setlength{\parskip}{1.2mm}

\newgeometry{left=2.5cm,right=2.5cm,top=2.5cm,bottom=2.5cm}
\begin{titlepage}
  \thispagestyle{empty}
  \setlength{\parindent}{0pt}
  \setlength{\parskip}{0pt}
  \fontsize{14}{18}\selectfont

  \begin{center}
    \textbf{ข้อเสนอโครงงานวิศวกรรมคอมพิวเตอร์}\\[4pt]
    \textbf{หลักสูตรวิศวกรรมศาสตรบัณฑิต สาขาวิชาวิศวกรรมคอมพิวเตอร์}\\[4pt]
    \textbf{ปีการศึกษา 2569}
  \end{center}
  \vspace{9mm}

  \textbf{หัวข้อโครงงาน} \par
  \hspace*{12mm} \thaidisstitleone
  \vspace{5mm}

  \textbf{สมาชิกกลุ่ม}\par
  \begingroup
  \fontsize{14}{18}\selectfont
  \setlength{\tabcolsep}{3pt}
  \renewcommand{\arraystretch}{1}
  \begin{tabularx}{\textwidth}{@{}r>{\RaggedRight\arraybackslash}Xl@{}}
    \hspace*{12mm}\dissstudentidone & \thaidissauthorone \\
    \dissstudentidtwo & \thaidissauthortwo \\
    \dissstudentidthree & \thaidissauthorthree \\
    \dissstudentidfour & \thaidissauthorfour \\
  \end{tabularx}
  \endgroup
  \par\vspace{5mm}

  \textbf{อาจารย์ที่ปรึกษา}\par
  \hspace*{12mm} \thaidissadvisor\par
  \vspace{5mm}

  \textbf{อาจารย์ที่ปรึกษาร่วม}\par
  \hspace*{12mm} \thaidisscoadvisor\par

  \vspace{48mm}

  \fontsize{13}{16}\selectfont By signing this, I hereby acknowledge that I have read the proposal and approved this Project\par
  \vspace{32mm}

  \begin{tabular*}{\textwidth}{@{\extracolsep{\fill}}>{\centering\arraybackslash}p{0.45\textwidth}>{\centering\arraybackslash}p{0.45\textwidth}@{}}
    & \hrulefill \\
    & \thaidissadvisor \\
    & อาจารย์ที่ปรึกษา \\
  \end{tabular*}

\end{titlepage}

%%%%%%%%%%%%%%%%%%%%% List of vocabs & terms %%%%%%%%%%%%%%%%%%%
%: Add your own terms and definitions below, following the example format.
\newcommand{\tableraglf}[1]{\RaggedRight{#1}}
\listofvocab
\begin{flushleft}
  \begin{xltabular}{\textwidth}{@{}>{\RaggedRight\arraybackslash}p{0.31\textwidth}@{\quad=\quad}>{\RaggedRight\arraybackslash}X@{}}
    UX/UI Design / การออกแบบประสบการณ์ผู้ใช้และส่วนต่อประสานผู้ใช้ & \tableraglf{กระบวนการออกแบบที่ครอบคลุมสองส่วนหลัก คือ UX (User Experience) ซึ่งเน้นการออกแบบภาพรวมของประสบการณ์ที่ผู้ใช้ได้รับขณะใช้งานผลิตภัณฑ์หรือบริการ เช่น ความง่ายในการใช้งาน ความสะดวก และความพึงพอใจ และ UI (User Interface) ซึ่งเน้นการออกแบบส่วนต่อประสานที่ผู้ใช้มองเห็นและโต้ตอบด้วยโดยตรง เช่น ปุ่มกด เมนู สี และรูปแบบตัวอักษร โดยทั้งสองส่วนทำงานร่วมกันเพื่อให้ผลิตภัณฑ์ดิจิทัล เช่น เว็บไซต์หรือแอปพลิเคชัน ใช้งานได้ง่าย มีประสิทธิภาพ และตอบสนองความต้องการของผู้ใช้ได้อย่างแท้จริง} \\

    UX/UI development life cycle (UX/UI DLC) & \tableraglf{วงจรการพัฒนางานออกแบบ UX/UI ซึ่งเป็นกระบวนการที่มีขั้นตอนต่อเนื่องกันตั้งแต่การเริ่มต้นจนถึงการนำไปใช้งานจริง โดยทั่วไปประกอบด้วยขั้นตอนหลัก เช่น การวิจัยและทำความเข้าใจผู้ใช้ (Research) การกำหนดปัญหาและความต้องการ (Define) การออกแบบแนวคิดและโครงร่าง (Ideate) การสร้างต้นแบบ (Prototype) และการทดสอบกับผู้ใช้จริง (Testing) ทั้งนี้ในแต่ละรอบของวงจรมักมีการนำผลตอบรับกลับมาปรับปรุงซ้ำ (Iteration) เพื่อให้ผลิตภัณฑ์ตอบโจทย์ผู้ใช้ได้ดียิ่งขึ้น} \\

    Usability Testing / การทดสอบการใช้งาน & \tableraglf{กระบวนการทดสอบการใช้งาน โดยให้ผู้เข้าร่วมทดสอบลองใช้งานระบบและหาจุดที่ใช้งานยากหรือติดขัด ซึ่งในขณะที่ทำการทดสอบ ผู้เข้าร่วมการทดสอบจำเป็นต้องคิดไปด้วยพูดไปด้วย (Thinking-Aloud) โดยนักออกแบบจะคอยสังเกตพฤติกรรมการใช้งานและรับฟังความเห็นของผู้เข้าร่วมทดสอบขณะทำภารกิจที่กำหนด} \\

    Comparative Usability Testing & \tableraglf{วิธีการทดสอบการใช้งานเชิงคุณภาพกับผู้เข้าร่วมทดสอบ โดยให้ผู้เข้าร่วมทดสอบทดลองใช้งานต้นแบบผลิตภัณฑ์สองรูปแบบที่มีเป้าประสงค์เดียวกันเพื่อเปรียบเทียบประสิทธิภาพการใช้งานและให้ข้อคิดเห็นหรือข้อเสนอแนะเกี่ยวกับรูปแบบหรือผลิตภัณฑ์นั้น ๆ พร้อมทั้งเลือกรูปแบบใดรูปแบบหนึ่งที่เห็นว่าเหมาะสมกับเป้าประสงค์นั้น ๆ} \\

    Contextual Feedback / ข้อเสนอแนะเชิงบริบท & \tableraglf{ข้อคิดเห็นหรือคำแนะนำที่กล่าวเฉพาะเจาะจงกับสิ่งใดสิ่งหนึ่ง เพื่อให้คำแนะนำมีความเกี่ยวข้องและนำไปใช้ได้ตรงจุดมากขึ้น} \\

    Wireframe / โครงร่างหน้าจอ & \tableraglf{โครงร่างพื้นฐานแสดงตำแหน่งองค์ประกอบบนหน้าจอ โดยยังไม่เน้นสีสันหรือรายละเอียดด้านภาพ ใช้สื่อสารการจัดวางในขั้นต้นของการออกแบบ} \\

    Prototype / ต้นแบบ & \tableraglf{ต้นแบบผลิตภัณฑ์ที่สร้างขึ้นเพื่อทดสอบแนวคิดก่อนพัฒนาจริง สามารถโต้ตอบได้ในระดับหนึ่ง แบ่งเป็นสามระดับตามความน้อย-มากของรายละเอียดในต้นแบบ คือ low, mid, high-fidelity prototype} \\

    Design System & \tableraglf{ชุดมาตรฐานการออกแบบที่รวบรวมองค์ประกอบ กฎเกณฑ์ และแนวทางปฏิบัติไว้ในที่เดียว เช่น สี ตัวอักษร ไอคอน และ component ต่าง ๆ เพื่อให้การออกแบบผลิตภัณฑ์มีความสอดคล้องและเป็นมาตรฐานเดียวกันทั้งระบบ} \\

    Insight / ข้อมูลเชิงลึก & \tableraglf{ความเข้าใจที่ได้จากการวิเคราะห์ข้อมูลหรือพฤติกรรมผู้ใช้ ซึ่งเปิดเผยปัญหาหรือความต้องการที่ซ่อนอยู่} \\

    User / ผู้ใช้ & \tableraglf{บุคคลที่มีปฏิสัมพันธ์โดยตรงกับผลิตภัณฑ์ ระบบ หรือบริการ เพื่อบรรลุเป้าหมายของตนเอง} \\

    Generative AI / ปัญญาประดิษฐ์แบบรู้สร้าง & \tableraglf{เทคโนโลยีปัญญาประดิษฐ์ที่สามารถสร้างเนื้อหาใหม่ เช่น ข้อความ ภาพ เสียง หรือโค้ดจากคำสั่ง (prompt)} \\

    Design Intent / ความต้องการทางการออกแบบ & \tableraglf{เจตนาหรือเป้าหมายเชิงการออกแบบที่นักออกแบบต้องการสื่อสาร ใช้เป็นแนวทางตัดสินใจในการออกแบบ} \\

    Component & \tableraglf{ส่วนประกอบย่อยที่นำกลับมาใช้ซ้ำได้ใน UI เช่น ปุ่มหรือช่องกรอกข้อมูล ช่วยให้การออกแบบสอดคล้องกัน} \\

    Two alternative UI & \tableraglf{UI ที่ใช้สำหรับการทำ comparative usability testing โดย UI ทั้งสองแบบต้องถูกทำขึ้นเพื่อเป้าประสงค์เดียวกัน แต่มีความแตกต่างกันในด้านการใช้งาน เช่น microcopy, layout} \\

    Human-in-the-loop & \tableraglf{แนวคิดที่มนุษย์เข้ามาใช้วิจารณญาณตรวจสอบและปรับปรุงผลลัพธ์จาก Generative AI เพื่อความถูกต้องเหมาะสม} \\

    Stakeholder / ผู้มีส่วนได้เสีย & \tableraglf{บุคคลหรือกลุ่มที่มีส่วนได้ส่วนเสียกับโครงการ เช่น ผู้บริหารหรือผู้ว่าจ้าง ซึ่งมีบทบาทในการตัดสินใจหรือสนับสนุนโครงการ} \\

   % Cognitive Overload & \tableraglf{ภาวะที่ได้รับข้อมูลมากเกินความสามารถในการประมวลผลของสมอง ส่งผลให้สับสนหรือตัดสินใจผิดพลาด} \\

    Requirement / ความต้องการ & \tableraglf{เงื่อนไขที่ผลิตภัณฑ์ต้องมีเพื่อตอบสนองความต้องการผู้ใช้หรือเป้าประสงค์ของโครงการ} \\

    Usability Tester / ผู้เข้าร่วมทดสอบ & \tableraglf{บุคคลที่เข้าร่วมทดสอบการใช้งานผลิตภัณฑ์ พร้อมให้ข้อคิดเห็นเพื่อนำไปปรับปรุงการออกแบบ} \\

    Persona Artifact & \tableraglf{เอกสารจำลองลักษณะผู้ใช้กลุ่มเป้าหมาย เช่น ชื่อ อายุ พฤติกรรม และเป้าหมาย เพื่อให้ทีมออกแบบเข้าใจผู้ใช้ชัดเจนขึ้น} \\
    %  & \tableraglf{ } \\

    %  & \tableraglf{ } \\

    %  & \tableraglf{ } \\

    %  & \tableraglf{ } \\

    %  & \tableraglf{ } \\
  \end{xltabular}
\end{flushleft}

\restoregeometry
